%% 1. CÀI ĐẶT CƠ BẢN & LỚP VĂN BẢN
\usepackage[utf8]{inputenc}
% \usepackage[T5]{fontenc} % Bỏ dòng này vì đã có T1 bên dưới, tránh xung đột
\usepackage[T1]{fontenc}
\usepackage[table,xcdraw]{xcolor}
\usepackage{tikz}
\usepackage{tcolorbox}

%% 2. BỐ CỤC TRANG & CANH LỀ (ĐÃ SỬA THEO YÊU CẦU)
% Theo TemplateCacMonDuAnVaThucTap_KhoaCNTT: A4, lề trên 3,5 - dưới 3 - trái 3,5 - phải 2 cm,
% số trang nằm ở giữa lề trên (header cách mép giấy 1,27 cm).
\usepackage[a4paper, top=3.5cm, bottom=3cm, left=3.5cm, right=2cm, headheight=16pt, headsep=1.65cm]{geometry}
\usepackage{fancyhdr}
\fancypagestyle{baocao}{%
  \fancyhf{}%
  \fancyhead[C]{\normalfont\fontsize{13pt}{15.6pt}\selectfont\thepage}%
  \renewcommand{\headrulewidth}{0pt}%
  \renewcommand{\footrulewidth}{0pt}%
}
\fancypagestyle{plain}{%
  \fancyhf{}%
  \fancyhead[C]{\normalfont\fontsize{13pt}{15.6pt}\selectfont\thepage}%
  \renewcommand{\headrulewidth}{0pt}%
  \renewcommand{\footrulewidth}{0pt}%
}
\pagestyle{baocao}
% Đã BỎ tùy chọn "twoside" để lề trang không bị lệch


%% 3. NGÔN NGỮ, FONT CHỮ & ĐỊNH DẠNG ĐOẠN VĂN (ĐÃ SỬA THEO YÊU CẦU)
\usepackage[vietnamese]{babel}
\usepackage{mathptmx} % Font Times New Roman
\usepackage{indentfirst}
% --- CÀI ĐẶT FONT, GIÃN DÒNG, THỤT LỀ MỚI ---
\renewcommand{\normalsize}{\fontsize{13}{15.6}\selectfont}
\linespread{1.5} % Giãn dòng 1.5
\setlength{\parindent}{1.27cm} % Thụt đầu dòng 1,27 cm (style "Nội dung văn bản" của template)
\setlength{\parskip}{0pt} % Template: không có khoảng cách trước/sau đoạn
\setlength{\emergencystretch}{2em}


%% 4. TÙY CHỈNH CÁC CẤP ĐỀ MỤC (SECTION, SUBSECTION,...)
\usepackage{titlesec}
\setcounter{secnumdepth}{4}

% Theo template khoa CNTT:
%   Heading 1 (Chương):          "CHƯƠNG N. TÊN CHƯƠNG", 16pt, đậm, canh trái
%   Heading 2 (Tiểu mục cấp 1):  14pt, đậm
%   Heading 3 (Tiểu mục cấp 2):  14pt, đậm nghiêng
%   Heading 4 (Tiểu mục cấp 3):  14pt, thường
%   Khoảng cách trước/sau mỗi đề mục 6pt, giãn dòng 1,5.
\titleformat{\section}[block]
  {\normalfont\fontsize{16pt}{19.2pt}\selectfont\bfseries\raggedright}
  {CHƯƠNG \thesection.}
  {0.4em}
  {}
\titlespacing*{\section}{0pt}{6pt}{6pt}

% Tiêu đề các trang mở đầu, danh mục, tài liệu tham khảo, phụ lục: 16pt, đậm, canh giữa
\titleformat{name=\section,numberless}[block]
  {\normalfont\fontsize{16pt}{19.2pt}\selectfont\bfseries\centering}
  {}
  {0em}
  {}
\titlespacing*{name=\section,numberless}{0pt}{6pt}{12pt}

\titleformat{\subsection}
  {\normalfont\fontsize{14pt}{16.8pt}\selectfont\bfseries\raggedright}
  {\thesubsection}
  {0.6em}
  {}
\titlespacing*{\subsection}{0pt}{6pt}{6pt}

\titleformat{\subsubsection}
  {\normalfont\fontsize{14pt}{16.8pt}\selectfont\bfseries\itshape\raggedright}
  {\thesubsubsection}
  {0.6em}
  {}
\titlespacing*{\subsubsection}{0pt}{6pt}{6pt}

\titleformat{\paragraph}
  {\normalfont\fontsize{14pt}{16.8pt}\selectfont\raggedright}
  {\theparagraph}
  {0.6em}
  {}
\titlespacing*{\paragraph}{0pt}{6pt}{6pt}


%% 5. HÌNH ẢNH, BẢNG BIỂU & CHÚ THÍCH (CAPTION)
\usepackage[nottoc]{tocbibind}
\usepackage{tocloft}
% Mục lục theo template: cấp 1 "CHƯƠNG N. TÊN" đậm, cấp 2 thường, cấp 3 nghiêng (thụt 0,25 cm),
% tất cả 13pt, có dấu chấm dẫn tới số trang.
\setlength{\cftbeforesecskip}{6pt}
\setlength{\cftbeforesubsecskip}{2pt}
\setlength{\cftbeforesubsubsecskip}{1pt}
\renewcommand{\cftsecpresnum}{CHƯƠNG }
\renewcommand{\cftsecaftersnum}{.}
\renewcommand{\cftsecleader}{\bfseries\cftdotfill{\cftdotsep}}
\cftsetpnumwidth{2.4em}
\cftsetrmarg{3em}
\setlength{\cftsubsecindent}{0pt}
\setlength{\cftsubsecnumwidth}{2.4em}
\setlength{\cftsubsubsecindent}{0.25cm}
\setlength{\cftsubsubsecnumwidth}{3.7em}
\renewcommand{\cftsubsubsecfont}{\itshape}
% Danh mục hình vẽ / bảng biểu: "Hình 2.1: Tên hình", "Bảng 4.1: Tên bảng"
\setlength{\cftfigindent}{0pt}
\renewcommand{\cftfigpresnum}{Hình }
\renewcommand{\cftfigaftersnum}{:}
\setlength{\cfttabindent}{0pt}
\renewcommand{\cfttabpresnum}{Bảng }
\renewcommand{\cfttabaftersnum}{:}
\AtBeginDocument{%
  \settowidth{\cftsecnumwidth}{\normalsize\bfseries CHƯƠNG 0.\ }%
  \settowidth{\cftfignumwidth}{\normalsize Hình 0.00:\ \ }%
  \settowidth{\cfttabnumwidth}{\normalsize Bảng 0.00:\ \ }%
}

% Căn giữa các tiêu đề danh mục tương tự tiêu đề section gốc
\renewcommand{\cfttoctitlefont}{\hspace*{\fill}\fontsize{16pt}{19.2pt}\selectfont\bfseries}
\renewcommand{\cftaftertoctitle}{\hspace*{\fill}}
\renewcommand{\cftloftitlefont}{\hspace*{\fill}\fontsize{16pt}{19.2pt}\selectfont\bfseries}
\renewcommand{\cftafterloftitle}{\hspace*{\fill}}
\renewcommand{\cftlottitlefont}{\hspace*{\fill}\fontsize{16pt}{19.2pt}\selectfont\bfseries}
\renewcommand{\cftafterlottitle}{\hspace*{\fill}}
\usepackage{graphicx}
\usepackage{float}
\usepackage{longtable}
\usepackage{tabularx}
\usepackage{multirow}
\usepackage[font=normal]{caption}
\usepackage{array}

% Chú thích theo template: "Hình 2.1: ...", "Bảng 4.1: ...", 13pt, canh giữa;
% chú thích bảng đặt phía trên bảng, chú thích hình đặt phía dưới hình.
\captionsetup{labelsep=colon, justification=centering, skip=6pt}
\captionsetup[table]{position=top}

\newcolumntype{s}{>{\hsize=.3\hsize}X}
\newcolumntype{y}{>{\hsize=.4\hsize}X}
\newcolumntype{d}{>{\hsize=.1\hsize}X}
\newcolumntype{a}{>{\hsize=1.1\hsize}X}
\newcolumntype{g}{>{\hsize=5\hsize}X}
% Căn chỉnh bảng: P{w} và X căn giữa theo cả chiều ngang và chiều dọc;
% L{w} và Y căn trái, giữa theo chiều dọc, dùng cho cột chứa đoạn văn dài.
\newcolumntype{P}[1]{>{\centering\arraybackslash}m{#1}}
\newcolumntype{L}[1]{>{\raggedright\arraybackslash}m{#1}}
\renewcommand{\tabularxcolumn}[1]{m{#1}}
\newcolumntype{Y}{>{\raggedright\arraybackslash}X}
\newcolumntype{Z}{>{\centering\arraybackslash}X}
% Giữ chữ trong các bảng ở mức dễ đọc, đồng nhất 12pt.
\usepackage{etoolbox}
\newcommand{\tabletextsize}{\fontsize{12pt}{14.4pt}\selectfont}
\AtBeginEnvironment{tabular}{\tabletextsize}
\AtBeginEnvironment{tabularx}{\tabletextsize}
\AtBeginEnvironment{longtable}{\tabletextsize}
% -----------------------


%% 6. KHỐI LỆNH (CODE LISTINGS) 
\usepackage{listings}
% Gói xcolor đã được nạp ở mục 5, không cần nạp lại

\lstdefinestyle{customcode}{
    basicstyle=\small\ttfamily,
    breaklines=true,
    numbers=left,
    numberstyle=\tiny\color{gray},
    keywordstyle=\color{blue!80!black}\bfseries,
    commentstyle=\color{green!50!black}\itshape,
    stringstyle=\color{red!70!black},
    showstringspaces=false,
    tabsize=2,
    captionpos=b,
    frame=single,
    rulecolor=\color{gray!40},
    xleftmargin=1.5em,
    framexleftmargin=1em,
    aboveskip=1.5em,
    belowskip=1.5em
}
\renewcommand{\lstlistingname}{Khối lệnh}


%% 7. TOÁN HỌC & CÁC MÔI TRƯỜNG ĐỊNH LÝ
\newtheorem{theorem}{Định lý}[section]
\newtheorem{defn}[theorem]{Định nghĩa}
\newtheorem{corollary}[theorem]{Hệ quả}
\newtheorem{lemma}[theorem]{Bổ đề}


%% 8. TÀI LIỆU THAM KHẢO
\usepackage[natbibapa]{apacite} 
\bibliographystyle{apacite}
\makeatletter
\def\APAC@apcfile{english.apc}
\makeatother

%% 9. CÁC GÓI BỔ SUNG KHÁC
\usepackage{wrapfig}
% Gói tikz đã được nạp ở mục 1, không cần nạp lại
\usetikzlibrary{calc,backgrounds}
\usepackage{soul}
\usepackage{colortbl}
\usepackage{forloop}
\usepackage{enumitem}

\definecolor{LightCyan}{rgb}{0.88,1,1}
\newcounter{loopcntr}
\newcommand{\rpt}[2][1]{\forloop{loopcntr}{0}{\value{loopcntr}<#1}{#2}}
\setitemize{itemsep=0pt,leftmargin=1.25cm,topsep = 3pt}


%% 10. ĐÁNH SỐ TỰ ĐỘNG & LIÊN KẾT 
\usepackage{chngcntr}
\usepackage[unicode, plainpages=false, pdfpagelabels=true, hidelinks]{hyperref}

% Cấu hình bộ đếm cho các môi trường
\counterwithin{figure}{section}
\counterwithin{table}{section}
\counterwithin{equation}{section}
% Các counter khác đã được article class và titlesec xử lý, không cần thêm


%% 11. ĐỊNH NGHĨA LẠI TÊN CÁC THÀNH PHẦN CỦA TÀI LIỆU
\AtBeginDocument{%
  \renewcommand{\contentsname}{MỤC LỤC}
  \renewcommand{\listfigurename}{DANH MỤC HÌNH VẼ}
  \renewcommand{\listtablename}{DANH MỤC BẢNG BIỂU}
  \renewcommand{\refname}{TÀI LIỆU THAM KHẢO}
  \normalsize % Áp dụng cỡ chữ mặc định cho toàn bộ tài liệu
}


%% 12. CÁC BIẾN TOÀN CỤC SỬ DỤNG TRONG BÁO CÁO
\newcommand{\khoa}{CÔNG NGHỆ THÔNG TIN}
\newcommand{\mon}{DỰ ÁN CÔNG NGHỆ THÔNG TIN}
\newcommand{\chuyennganh}{MẠNG MÁY TÍNH VÀ TRUYỀN THÔNG DỮ LIỆU}
\newcommand{\de}{XÂY DỰNG MẠNG CAMPUS NETWORK SỬ DỤNG SDN VÀ SD-WAN}
\newcommand{\gvhd}{ThS. LÊ VIẾT THANH}
\newcommand{\svmot}{TRẦN HỮU NHÂN - 52300235}
\newcommand{\svhai}{NGUYỄN NHẬT HÀO - 52300198}
\newcommand{\lop}{23050401}
\newcommand{\khoahoc}{27}
\newcommand{\nam}{2026}
